\subsection{The finite norm-torus module and its local graphs}

\begin{statement}{A.109}{Proved}
{Actual finite norm-one classes, approximation, and the auxiliary local condition}
Let $a\in\Q^\times$ have $v_3(a)=1$, let $N=3^n$ with
$n\ge1$, and put $K=\Q(\alpha)$, $\alpha^N=a$.
For $r\ge1$ write
\[
 R_r=\Z/3^r\Z,\qquad d=\max(n,r),\qquad
 L=\Q(\alpha,\mu_{3^d}),\qquad
 G=C_N\rtimes U_d,\quad U_d=(\Z/3^d\Z)^\times.
 \tag{A.109.1}
\]
The displayed group is the actual group $\operatorname{Gal}(L/\Q)$.
Let $H=U_d$ be the stabilizer of $\alpha$, and define
\[
 P_r=R_r[G/H],\quad I_r=\ker(P_r\xrightarrow{\sum}R_r),
 \quad J_r=P_r/R_r\mathbf1,
 \quad T=\ker\bigl(\operatorname{Res}_{K/\Q}\mathbf G_m
                         \xrightarrow{N}\mathbf G_m\bigr).
 \tag{A.109.2}
\]
There are canonical identifications of finite $G_\Q$-modules
\[
 M_r:=T[3^r]=I_r\otimes\mu_{3^r},\qquad
 M_r^\vee:=\operatorname{Hom}(M_r,\mu_{3^r})=J_r.
 \tag{A.109.3}
\]
Both $H^1(G,M_r)$ and $H^1(G,J_r)$ vanish.
For either $C=M_r$ or $C=J_r$ and every finite set $S$ of places,
the actual localization map
\[
 H^1(\Q,C)\longrightarrow\prod_{v\in S}H^1(\Q_v,C)
 \tag{A.109.4}
\]
is onto.  Every nonzero global class has nonzero restriction to $G_L$
and is detected at infinitely many primes splitting completely in $L$.
This last assertion allows exclusion of a fixed finite set of primes.

The classes represented by actual global norm-one elements are exactly
\[
 \operatorname{im}\bigl(T(\Q)/T(\Q)^{3^r}\to H^1(\Q,M_r)\bigr)
 =\left\{c:\operatorname{loc}_v(c)\in
       \operatorname{im}\bigl(T(\Q_v)/T(\Q_v)^{3^r}\bigr)
          \text{ for every }v\right\}.
 \tag{A.109.5}
\]

At a prime $q\ne3$ splitting completely in $L$, choose compatible
local Kummer coordinates.  Then
\[
 H^1(\Q_q,M_r)=I_r\oplus I_r,\qquad
 H^1(\Q_q,J_r)=J_r\oplus J_r,
 \tag{A.109.6}
\]
where $(t,c)$ records tame inertia and Frobenius, respectively.
For $z\in R_r$ the graph
$\mathcal L_z=\{(t,zt):t\in I_r\}$ consists exactly of the
classes whose local image is generated by its tame inertia image
and whose chosen Frobenius equals $z$ times that inertia generator.
Its Tate orthogonal is the graph $\{(u,zu):u\in J_r\}$.
At precision three, every ramified class with cyclic local image lies
on one of these graphs.  The unramified classes form the separate
subspace $\{(0,c)\}$.

Here is the direct map to the next norm equation.  Suppose additionally
that $q$ splits the next pure-radical field, so all conjugates
$\alpha_i$ are cubes in $\Q_q$, and that $q$ is prime to the target
rational norm $b$, with $b$ a local cube.  Write the local coordinates
of an actual target $\beta$ of norm $b$ as
$t_i=v_q(\beta_i)\bmod3$ and
$c_i=[q^{-v_q(\beta_i)}\beta_i]$ in the unit cubic quotient.
For $Q=q^kQ_0$, $R$ a $q$-unit, and $3\nmid ek$, the simultaneous
local conditions
\[
 \beta_iR^e\in
 N_{\Q_q((Q^e\alpha_i)^{1/3})/\Q_q}
           \Q_q((Q^e\alpha_i)^{1/3})^\times
 \quad\text{for every }i
 \tag{A.109.7}
\]
are exactly
\[
 c=z t+w\mathbf1,\qquad
       z=[Q_0]/k,\qquad w=-e[R]\quad\text{in }\F_3.
 \tag{A.109.8}
\]
Thus a locally cubic $R$ gives the graph condition itself.
No assertion of unrestricted approximation with (A.109.8) at all
new ramification is made; A.110 gives an explicit obstruction to it.
\end{statement}

\begin{proof}
An odd cube root of $a$ cannot lie in an abelian extension of $\Q$:
its field has degree three by the three-adic valuation and is not
Galois.  Consequently $a$ has full $3^n$-Kummer order over the
cyclotomic field in (A.109.1).  Its radical automorphism and the
automorphisms fixing the chosen $\alpha$ give the stated split group.
On the geometric torus, a coordinate is attached to each embedding
$K\hookrightarrow\overline\Q$, and the norm is their product.
Taking $3^r$-torsion gives the literal exact sequence
\[
 0\longrightarrow M_r\longrightarrow P_r\otimes\mu_{3^r}
      \xrightarrow{\sum}\mu_{3^r}\longrightarrow0.
 \tag{A.109.9}
\]
The coordinate dot product pairs $I_r$ perfectly with
$P_r/R_r\mathbf1$; this proves both identifications in (A.109.3).

\emph{The two finite-group vanishings.}
The central involution $-1\in H$ acts on $\mu_{3^r}$ by $-1$.
Inner action on group cohomology is the identity, whereas this
coefficient action is minus the identity.  Since two is invertible,
\[
 H^i(H,\mu_{3^r})=0\quad(i\ge0),\qquad
 H^i(G,P_r\otimes\mu_{3^r})=0\quad(i\ge0)
 \tag{A.109.10}
\]
by Shapiro.  Also $H^0(G,\mu_{3^r})=0$.
The beginning of the exact sequence of (A.109.9) proves
$H^1(G,M_r)=0$.

For $J_r$ use $0\to R_r\to P_r\to J_r\to0$.
The map $H^i(G,R_r)\to H^i(G,P_r)=H^i(H,R_r)$ is restriction.
It is an isomorphism for $i=1,2$.
Indeed $G^{\rm ab}=H$, because the involution inverts $C_N$,
so the assertion for $i=1$ follows directly.
In the Hochschild--Serre sequence for $C_N\triangleleft G$, the groups
$H^1(C_N,R_r)=R_r[N]$ and $H^2(C_N,R_r)=R_r/NR_r$
are negated by that involution.
Their $H$-cohomology is therefore zero in every degree.
All terms with positive $C_N$-degree contributing to total degree
at most two vanish.  The split inclusion $H\hookrightarrow G$
identifies the remaining row with $H^i(H,R_r)$.
The long exact sequence now gives $H^1(G,J_r)=0$.

Neukirch's finite-module approximation theorem, in the precise form
\cite[Theorem 2.1(c)]{BIR_Neftin2013}, proves (A.109.4).
If the splitting group of the Cartier dual is a quotient of $G$,
its $H^1$ injects by inflation into the vanishing group just computed;
thus the theorem applies to the actual minimal splitting field as well.
Inflation--restriction proves that a nonzero class restricts
nontrivially to $G_L$.  On $G_L$ it is an ordinary homomorphism to
the finite module.  Choose an element on which it is nonzero and
apply Chebotarev to the actual finite Galois extension cut out by
the homomorphism $g\mapsto(c(g),\bar g)$ into $C\rtimes G$.
Its Frobenius class projects trivially to $G$ and has nonzero module
component; conjugation preserves nonvanishing.
This proves the detection assertion with its stated scope.

\emph{The exact norm-one comparison.}
The Kummer sequence for $T$ gives
\[
 \begin{gathered}
 0\longrightarrow T(\Q)/T(\Q)^{3^r}
 \longrightarrow H^1(\Q,M_r)
 \longrightarrow H^1(\Q,T)[3^r]\longrightarrow0,\\
 H^1(\Q,T)=\Q^\times/N_{K/\Q}K^\times.
 \end{gathered}
 \tag{A.109.11}
\]
There is the identical local sequence at every place.
The Hasse norm theorem for this actual pure-radical holomorph,
proved and sourced in A.77, says that the kernel of localization
on $H^1(\Q,T)$ is zero.
Applying these maps to (A.109.11) proves (A.109.5).
In particular (A.109.5) retains the norm-one condition, which is
stronger than merely asking that the norm of a Kummer representative
be a cube.

\emph{Local coordinates and the literal norm condition.}
At the specified $q$, all coefficient modules are trivial and
$\Q_q^\times/(\Q_q^\times)^{3^r}$ has its uniformizer and unit
cyclic factors of order $3^r$.
Choose tame inertia $\tau$ and a Frobenius $\Phi$ so that these
Kummer factors evaluate as $(1,0)$ and $(0,1)$ on $(\tau,\Phi)$.
Thus a cocycle has values $(t,c)$ as asserted.
Its local image is generated by $t,c$, and its inertia image by $t$.
This proves the graph assertions, including the distinction when
$r>1$ between cyclic image and generation by inertia.
With the unit generator oriented by local reciprocity, local Tate
duality is the perfect pairing
\[
 \langle(t,c),(u,v)\rangle_q
                 =\frac{t\cdot v-c\cdot u}{3^r}
                 \quad\text{in }\Q/\Z.
 \tag{A.109.12}
\]
The dot product $I_r\times J_r\to R_r$ is perfect, so substitution
gives exactly the stated orthogonal graph.

Finally the tame cubic norm formula proved in A.103, applied to
$q^{ek}Q_0^e\alpha_i$ and $\beta_iR^e$, is
\[
        t_i e[Q_0]-ek(c_i+e[R])=0\quad\text{in }\F_3.
 \tag{A.109.13}
\]
The class of $\alpha_i$ is zero by the specified deeper splitting.
Dividing by $ek$ gives (A.109.8), and proves the equivalence with
(A.109.7).  The assumptions on $b$ imply
$\sum_i t_i=\sum_i c_i=0$.
Also $\mathbf1\in I_1$ because $3\mid N$, so the scalar direction
in (A.109.8) is a genuine direction in the stated local module.
\end{proof}

The retained script \texttt{norm\_torus\_selmer\_certificate.py}
computes the cocycles and coboundaries from the actual two-generator
holomorph presentation using integer Smith normal forms.
Its output verifies both vanishings in fourteen finite cases,
including precisions three, nine, and twenty-seven.
The proof above, rather than those finite cases, establishes all
$n,r$.
